%% Font size %%
\documentclass[11pt]{article}

%% Load the custom package
\usepackage{Mathdoc}

%% Numéro de séquence %% Titre de la séquence %%
\renewcommand{\centerhead}{Fractions - Éval 2 : Droites graduées}

%% Spacing commands %%
\renewcommand{\baselinestretch}{1} \setlength{\parindent}{0pt}

%% Droite graduée vide, unité #1 sur \droitex cm, #2 = nombre max (entier), #3 = graduations par unité
\newcommand{\droitex}{2.2cm}
\newcommand{\droitevide}[2]{%
\begin{tikzpicture}[baseline=-0.5ex, x=\droitex]
\useasboundingbox (-0.3,-0.5) rectangle (#1+0.4,0.7);
\draw[thick,-latex] (-0.3,0) -- (#1+0.4,0);
\foreach \i in {0,...,#1}{
\draw[thick] (\i,-4pt) -- (\i,4pt);
\node[below] at (\i,-4pt) {\i};
}
\foreach \i in {0,...,\numexpr#1*#2\relax}{
\pgfmathsetmacro{\x}{\i/#2}
\pgfmathtruncatemacro{\reste}{mod(\i,#2)}
\ifnum\reste=0\relax\else
\draw (\x,-2pt) -- (\x,2pt);
\fi
}
\end{tikzpicture}%
}

%% Droite graduée avec points placés, #1 = max, #2 = graduations par unité, #3 = liste "x/Lettre,..."
\newcommand{\droitepoints}[3]{%
\begin{tikzpicture}[baseline=-0.5ex, x=\droitex]
\useasboundingbox (-0.3,-0.5) rectangle (#1+0.4,0.7);
\draw[thick,-latex] (-0.3,0) -- (#1+0.4,0);
\foreach \i in {0,...,#1}{
\draw[thick] (\i,-4pt) -- (\i,4pt);
\node[below] at (\i,-4pt) {\i};
}
\foreach \i in {0,...,\numexpr#1*#2\relax}{
\pgfmathsetmacro{\x}{\i/#2}
\pgfmathtruncatemacro{\reste}{mod(\i,#2)}
\ifnum\reste=0\relax\else
\draw (\x,-2pt) -- (\x,2pt);
\fi
}
\foreach \px/\lettre in {#3}{
\filldraw (\px,0) circle (2pt);
\node[above] at (\px,4pt) {\lettre};
}
\end{tikzpicture}%
}

\begin{document}

\phantom{0}
\vspace{-.5cm}

\entetedevoirs{20}

\begin{center}
\duree{30 minutes}
\total{20 points}
\soin
\end{center}

\begin{exercicedevoir}[2,5][Placer des demis]
Place les points suivants sur la droite graduée : $A = \dfrac{1}{2}$ ; $B = \dfrac{5}{2}$ ; $C = 3$ ; $D = \dfrac{7}{2}$ ; $E = \dfrac{3}{2}$.
\begin{center}
\droitevide{4}{2}
\end{center}
\end{exercicedevoir}

\begin{exercicedevoir}[2,5][Placer des tiers]
Place les points suivants sur la droite graduée : $A = \dfrac{4}{3}$ ; $B = 2$ ; $C = \dfrac{8}{3}$ ; $D = \dfrac{5}{3}$ ; $E = \dfrac{11}{3}$.
\begin{center}
\droitevide{4}{3}
\end{center}
\end{exercicedevoir}

\begin{exercicedevoir}[2,5][Placer des quarts]
Place les points suivants sur la droite graduée : $A = \dfrac{3}{4}$ ; $B = \dfrac{9}{4}$ ; $C = 3$ ; $D = \dfrac{13}{4}$ ; $E = \dfrac{7}{4}$.
\begin{center}
\droitevide{4}{4}
\end{center}
\end{exercicedevoir}

\begin{exercicedevoir}[2,5][Placer des sixièmes]
Place les points suivants sur la droite graduée : $A = \dfrac{5}{6}$ ; $B = \dfrac{8}{6}$ ; $C = 2$ ; $D = \dfrac{13}{6}$ ; $E = \dfrac{16}{6}$.
\begin{center}
\droitevide{3}{6}
\end{center}
\end{exercicedevoir}

\newpage
\phantom{0}

\begin{exercicedevoir}[2,5][Lire des demis]
Lis l'abscisse de chaque point, puis complète.
\begin{center}
\droitepoints{5}{2}{3.5/A,0.5/B,4.5/C,1.5/D,2.5/E}
\end{center}
$A = \dotfill$ \quad $B = \dotfill$ \quad $C = \dotfill$ \quad $D = \dotfill$ \quad $E = \dotfill$
\end{exercicedevoir}

\begin{exercicedevoir}[2,5][Lire des tiers]
Lis l'abscisse de chaque point, puis complète.
\begin{center}
\droitepoints{5}{3}{1.667/A,4.667/B,0.333/C,3/D,3.667/E}
\end{center}
$A = \dotfill$ \quad $B = \dotfill$ \quad $C = \dotfill$ \quad $D = \dotfill$ \quad $E = \dotfill$
\end{exercicedevoir}

\begin{exercicedevoir}[2,5][Lire des quarts]
Lis l'abscisse de chaque point, puis complète.
\begin{center}
\droitepoints{5}{4}{1.75/A,3.75/B,0.75/C,4.75/D,2.75/E}
\end{center}
$A = \dotfill$ \quad $B = \dotfill$ \quad $C = \dotfill$ \quad $D = \dotfill$ \quad $E = \dotfill$
\end{exercicedevoir}

\begin{exercicedevoir}[2,5][Lire des sixièmes]
Lis l'abscisse de chaque point, puis complète.
\begin{center}
\droitepoints{4}{6}{2.833/A,0.5/B,2/C,3.667/D,1.333/E}
\end{center}
$A = \dotfill$ \quad $B = \dotfill$ \quad $C = \dotfill$ \quad $D = \dotfill$ \quad $E = \dotfill$
\end{exercicedevoir}

\end{document}
